\documentclass[../main.tex]{subfiles}
\begin{document}

\section{Overview}

\noindent This appendix provides the concrete commands, key functions, and reproduction steps for the edge deployment pipeline described in Chapter~6. All scripts referenced here are archived in the \texttt{deployment/} directory of the project repository.

\section{Prerequisites}

\begin{itemize}[itemsep=0.4em, topsep=0.4em]
    \item Host machine: macOS or Linux with Python 3.10+.
    \item Packages: \texttt{torch}, \texttt{pytorch-lightning}, \texttt{hydra-core}, \texttt{peft}, \texttt{onnx}, \texttt{onnxruntime}, \texttt{qai-hub}.
    \item API token: Qualcomm AI Hub (\texttt{qai-hub configure --api\_token <TOKEN>}).
    \item Target device: Qualcomm RB3 Gen2 Vision Development Kit with Ubuntu 24.04 aarch64, accessible via SSH or Cockpit.
\end{itemize}

\section{Step 0: Checkpoint Analysis}

\noindent The \texttt{analyze\_checkpoint.py} script performs a dry-run memory projection before committing to export. It reports per-module parameter counts (FP32/FP16/INT8 sizes), total weights size, estimated activation memory (0.5$\times$ weights heuristic), and a verdict against the 4\,GB device budget.

% --- CODE BLOCK 1 ---
% Sử dụng basicstyle=\ttfamily\scriptsize để giảm cỡ chữ
% breaklines=true và breakatwhitespace=false để ép xuống dòng các dòng dài
\begin{lstlisting}[
    language=bash,
    caption={Checkpoint analysis command},
    label={lst:analyze_checkpoint},
    basicstyle=\ttfamily\scriptsize,
    breaklines=true,
    breakatwhitespace=false,
    xleftmargin=1.5em,
    xrightmargin=0.5em
]
python deployment/scripts/analyze_checkpoint.py \
    --ckpt epoch=56-val_score=52.28.ckpt
\end{lstlisting}

\section{Step 1: LoRA Merge and FP16 Export}

\noindent The \texttt{lora\_fp16/export.py} script merges LoRA adapters into the base weights and converts floating-point tensors to FP16.

% --- CODE BLOCK 2 ---
\begin{lstlisting}[
    language=bash,
    caption={LoRA merge + FP16 export command},
    label={lst:lora_export_cmd},
    basicstyle=\ttfamily\scriptsize,
    breaklines=true,
    breakatwhitespace=false,
    xleftmargin=1.5em,
    xrightmargin=0.5em
]
python deployment/scripts/lora_fp16/export.py \
    --ckpt epoch=56-val_score=52.28.ckpt \
    --output-dir exported_model
\end{lstlisting}

\noindent The core merge operation uses PEFT's built-in method, followed by tensor deduplication and FP16 conversion to preserve the SigLIP shared-weight structure:

% --- CODE BLOCK 3 ---
\begin{lstlisting}[
    language=Python,
    caption={LoRA merge and FP16 conversion with tensor deduplication},
    label={lst:lora_merge_fp16},
    basicstyle=\ttfamily\scriptsize,
    breaklines=true,
    breakatwhitespace=false,
    xleftmargin=1.5em,
    xrightmargin=0.5em
]
from peft import PeftModel

# Reconstruct the Lightning model with LoRA-wrapped backbone
model = LitTBPS(config=cfg)
model.load_state_dict(ckpt["state_dict"], strict=False)

# Merge adapters into base weights: W_merged = W_0 + alpha*(BA)/r
if isinstance(model.model.backbone, PeftModel):
    model.model.backbone = model.model.backbone.merge_and_unload()

# Deduplicate shared tensors via data_ptr() before FP16 conversion
seen = {}
fp16_state = {}
for key, tensor in model.state_dict().items():
    ptr = tensor.data_ptr()
    if ptr in seen:
        # Reuse already-converted tensor to preserve shared-weight structure
        fp16_state[key] = seen[ptr]
    else:
        # Convert floating-point tensors; keep integer tensors as-is
        converted = tensor.half() if tensor.is_floating_point() else tensor
        seen[ptr] = converted
        fp16_state[key] = converted

torch.save(fp16_state, "exported_model/model_fp16.pt")
\end{lstlisting}

\noindent Outputs:
\begin{itemize}[itemsep=0.3em, topsep=0.3em]
    \item \texttt{exported\_model/model\_fp32.pt} (1418.5\,MB on disk)
    \item \texttt{exported\_model/model\_fp16.pt} (709.3\,MB on disk)
    \item \texttt{exported\_model/config.yaml} (resolved Hydra config)
\end{itemize}

\section{Step 2: ONNX Conversion}

\noindent The \texttt{onnx/export.py} script converts the FP16 state dict into two separate ONNX encoders (vision and text), each exported as a directory containing both the graph protobuf and the external weights file.

% --- CODE BLOCK 4 ---
\begin{lstlisting}[
    language=bash,
    caption={ONNX conversion command},
    label={lst:onnx_export_cmd},
    basicstyle=\ttfamily\scriptsize,
    breaklines=true,
    breakatwhitespace=false,
    xleftmargin=1.5em,
    xrightmargin=0.5em
]
python deployment/scripts/onnx/export.py \
    --model-dir exported_model \
    --precision fp32
\end{lstlisting}

\noindent The core export uses \texttt{torch.onnx.export} with opset 18 and a dynamic batch axis; the vision tower wrapper is shown below:

% --- CODE BLOCK 5 ---
\begin{lstlisting}[
    language=Python,
    caption={ONNX export of the vision encoder},
    label={lst:onnx_vision},
    basicstyle=\ttfamily\scriptsize,
    breaklines=true,
    breakatwhitespace=false,
    xleftmargin=1.5em,
    xrightmargin=0.5em
]
class VisionWrapper(nn.Module):
    def __init__(self, model):
        super().__init__()
        self.model = model

    def forward(self, image: torch.Tensor) -> torch.Tensor:
        # (batch, 3, 256, 256) -> (batch, 768)
        return self.model.encode_image(image)

dummy = torch.randn(1, 3, 256, 256)
torch.onnx.export(
    VisionWrapper(model).eval(),
    dummy,
    "exported_model/vision_onnx/vision_encoder.onnx",
    opset_version=18,
    input_names=["image"],
    output_names=["image_embed"],
    dynamic_axes={
        "image": {0: "batch"},
        "image_embed": {0: "batch"},
    },
)
\end{lstlisting}

\noindent Outputs (each directory contains both graph and external weights, sizes measured on disk):
\begin{itemize}[itemsep=0.3em, topsep=0.3em]
    \item \texttt{exported\_model/vision\_onnx/vision\_encoder.onnx} (1.44\,MB)
    \item \texttt{exported\_model/vision\_onnx/vision\_encoder.onnx.data} (354.6\,MB)
    \item \texttt{exported\_model/text\_onnx/text\_encoder.onnx} (1.36\,MB)
    \item \texttt{exported\_model/text\_onnx/text\_encoder.onnx.data} (1059.3\,MB)
\end{itemize}

\section{Step 3: QNN Compilation via Qualcomm AI Hub}

\noindent Submit each encoder \textbf{directory} (not the individual \texttt{.onnx} file) as a compile job:

% --- CODE BLOCK 6 ---
\begin{lstlisting}[
    language=bash,
    caption={QNN compile jobs via Qualcomm AI Hub},
    label={lst:qnn_compile},
    basicstyle=\ttfamily\scriptsize,
    breaklines=true,
    breakatwhitespace=false,
    xleftmargin=1.5em,
    xrightmargin=0.5em
]
qai-hub submit-compile-job \
    --model exported_model/vision_onnx/ \
    --device "Dragonwing RB3 Gen 2 Vision Kit" \
    --compile_options "--target_runtime qnn_context_binary" \
    --name "mSigLIP-vision" --wait

qai-hub submit-compile-job \
    --model exported_model/text_onnx/ \
    --device "Dragonwing RB3 Gen 2 Vision Kit" \
    --compile_options "--target_runtime qnn_context_binary" \
    --name "mSigLIP-text" --wait
\end{lstlisting}

\noindent Target runtime options:
\begin{itemize}[itemsep=0.3em, topsep=0.3em]
    \item \texttt{qnn\_context\_binary} (recommended, pre-compiled for direct DSP execution)
    \item \texttt{qnn\_dlc} (legacy DLC format)
    \item \texttt{onnx} (fallback to ONNX Runtime CPU)
\end{itemize}

\noindent \textbf{Common pitfall:} Passing \texttt{--model exported\_model/vision\_onnx/vision\_encoder.onnx} (single file) instead of the directory fails with a cryptic shape-mismatch error because the external \texttt{.onnx.data} weights are not uploaded. Always pass the directory.

\section{Step 4: On-Device Inference}

\noindent On the RB3 Gen2 device, run inference via the QNN runtime with \texttt{qnn-net-run}, or fall back to ONNX Runtime CPU for intermediate validation.

% --- CODE BLOCK 7 ---
\begin{lstlisting}[
    language=bash,
    caption={On-device inference via QNN},
    label={lst:ondevice_inference},
    basicstyle=\ttfamily\scriptsize,
    breaklines=true,
    breakatwhitespace=false,
    xleftmargin=1.5em,
    xrightmargin=0.5em
]
# Transfer compiled binaries to device
scp vision.bin text.bin ubuntu@rb3:/home/ubuntu/msiglip/

# SSH in and run
ssh ubuntu@rb3
qnn-net-run --backend libQnnHtp.so \
    --model vision.bin \
    --input_list vision_inputs.txt \
    --output_dir vision_out/
\end{lstlisting}

\noindent A CPU fallback via ONNX Runtime (for regression testing against the PyTorch baseline):

% --- CODE BLOCK 8 ---
\begin{lstlisting}[
    language=bash,
    caption={ONNX Runtime CPU inference test},
    label={lst:onnx_inference_test},
    basicstyle=\ttfamily\scriptsize,
    breaklines=true,
    breakatwhitespace=false,
    xleftmargin=1.5em,
    xrightmargin=0.5em
]
python deployment/scripts/inference_test.py \
    --model-dir exported_model \
    --backend onnx \
    --dataset-root /path/to/VN3K
\end{lstlisting}

\section{Hardware Profiling (Proxy Models)}

\noindent The \texttt{hardware\_profiling/} suite benchmarks standard proxy models (MobileNetV2, ResNet18, EfficientNet-B0) on the target device. Output logs are written to \texttt{deployment/logs/} with a timestamped prefix for reproducibility.

% --- CODE BLOCK 9 ---
\begin{lstlisting}[
    language=bash,
    caption={Hardware profiling suite execution},
    label={lst:hardware_profiling},
    basicstyle=\ttfamily\scriptsize,
    breaklines=true,
    breakatwhitespace=false,
    xleftmargin=1.5em,
    xrightmargin=0.5em
]
cd deployment/hardware_profiling
bash install_deps.sh       # Install Python dependencies on device
bash collect_sysinfo.sh    # Dump CPU, RAM, and SoC information
python benchmark.py        # PyTorch CPU vs ONNX Runtime
                           # on MobileNetV2, ResNet18, EfficientNet-B0
\end{lstlisting}

\section{Reproducibility Checklist}

\begin{itemize}[itemsep=0.4em, topsep=0.4em]
    \item[$\square$] Checkpoint analyzed; FP16 verdict $=$ comfortable.
    \item[$\square$] \texttt{model\_fp16.pt} produced and validated via reload and forward equivalence.
    \item[$\square$] \texttt{vision\_onnx/} and \texttt{text\_onnx/} directories produced; each contains both \texttt{.onnx} and \texttt{.onnx.data}.
    \item[$\square$] AI Hub compile jobs submitted with directory (not file) model references.
    \item[$\square$] Compiled \texttt{.bin} files transferred to RB3 Gen2.
    \item[$\square$] \texttt{qnn-net-run} output matches ONNX Runtime CPU output within tolerance.
    \item[$\square$] Retrieval metrics (Rank-1, mAP) on VN3K test set preserved within 1\% absolute vs.\ PyTorch FP32.
\end{itemize}

\end{document}
